\section{Nested triangle insertions}
\label{sec:nested-triangle}

The three graphs L2E3G2, L2E3G3, L2E3G5 have one logarithmically divergent
proper triangle. They use the same massive Euclidean zero-momentum BPH
conditions as the rest of the atlas. The standard kernel below distinguishes
the physical leg $p_1$ attached to that triangle; the other two legs may be
exchanged. For the catalog graphs its arguments are, respectively,
$(s_1,s_2,s_3)$, $(s_2,s_1,s_3)$, and $(s_3,s_1,s_2)$.

\subsection{Standard routing and the four recursive terms}

All momenta are incoming, $p_1+p_2+p_3=0$. Set
\[
 a(k)=D(k)D(k+p_2)D(k-p_1),\qquad
 h(k,q)=D(q)D(q+k)D(q+k-p_1).
\]
The subgraph Taylor operator sets its independent external momenta
$(k,p_1)$ to zero at fixed $q$, so its counterterm is $-D(q)^3$.
The overall operator sets the physical $(p_1,p_2)$ to zero at fixed
$(k,q)$. Both orders are zero. Thus
\begin{align*}
 RI={}&a(k)\{h(k,q)-D(q)^3\}\\
 &-D(k)^3\{D(q)D(q+k)^2-D(q)^3\}.
\end{align*}
These four terms are the direct recursive BPH expression, not a change to
an MS finite prescription. In a standard fixed-label realization the
propagator-order edges are
$(1,3),(1,2),(2,4),(3,4),(0,3),(0,4)$, and external labels $1,2,3$ attach
to vertices $0,1,2$. This identifies the distinguished physical leg without
relying on a symmetric momentum example.

\subsection{Six common Schwinger labels and Gaussian integration}

At positive cutoff replace every propagator by $D_\eta$ before Taylor
subtraction. The four terms retain the same six original Schwinger labels
and common domain $\alpha_e\ge\eta$. Split the parameters into outer sum
$A$ and inner sum $C$, and later impose $A+C=1$. Let the outer fractions be
$(a,b,c)$ and the inner fractions $(u,v,w)$, with each triple nonnegative
and summing to one, in the propagator orders displayed above. Write
\begin{align*}
 H&=u(1-u),& B&=uvw,& D_*&=A+CH,\\
 U_{\rm raw}&=CD_*,& U_{\rm ct}&=CA,\\
 k_0&=cp_1-bp_2,&
 W&=ab s_2+ac s_1+bc s_3,\\
 G&=uv|k_0|^2+uw|k_0-p_1|^2+vw s_1,\\
 \delta&=\frac{C(AG+CBs_1)}{D_*}.
\end{align*}
The notation $D_*$ is a parameter denominator, not a propagator.
Completing the $q$ square first and then the $k$ square gives external
variance $AW+\delta$ for the raw graph and $AW$ for its subgraph
counterterm. In particular
\[
 F_{\rm raw}=U_{\rm raw}\{M(A+C)+AW+\delta\}.
\]
These identities hold before normalizing $A+C$, and may be verified by
independent spanning-tree and two-forest polynomials.

For the grouped parameters, the original six-edge measure is
\[
 A^2C^2\,dA\,dC\,da\,db\,du\,dv,
\]
with $c=1-a-b$, $w=1-u-v$. On the common projective simplex the remaining
factor is $A^2C^2\,dA$ times the two fraction-simplex measures. The radial
six-parameter measure is $\rho^5d\rho$, and the two six-dimensional
Gaussian factors supply $\rho^{-6}$. Consequently the overall subtraction
produces the logarithmic kernel
\[
 L(z)=\log(1+z/M).
\]
At positive cutoff it is instead
\[
 L_T(z)=\int_T^\infty\frac{d\rho}{\rho}
           e^{-M\rho}(1-e^{-z\rho}),\qquad
 T=\eta/\min_e x_e,
\]
and
\[
 0\le L_T(z)\le L(z),\qquad
 0\le L_T'(z)\le\frac1{M+z}.
\]
The same $T$ occurs in both paired terms. The candidate finite integral is
\begin{align*}
 J_{\mathrm{tr}}
 &=\int_0^1dA\int_{\Delta_o}da\,db\int_{\Delta_i}du\,dv\,\mathcal F,\\
 \mathcal F
 &=\frac{A^2}{C}
       \left\{\frac{L(AW)}{A^3}
                    -\frac{L(AW+\delta)}{D_*^3}\right\},\qquad C=1-A.
\end{align*}

\subsection{Exact pairing at the subdivergent face}

The two terms must be paired before $C\to0$ or the common regulator is
removed. Define $\ell(z)=\log(1+z)/z$, with $\ell(0)=1$. Then
\begin{align*}
 \mathcal F={}&
 \frac{H(D_*^2+AD_*+A^2)}{D_*^3}
       \frac WM\ell(AW/M)\\
 &-\frac{A^2(AG+CBs_1)}{D_*^4(M+AW)}
                       \ell\left(\frac\delta{M+AW}\right).
\end{align*}
There is no $1/C$ left in this stable expression. For estimates write it
as $P_+-P_-$, where the following two quantities are nonnegative:
\begin{align*}
 P_+&=\frac{A^2}{C}(A^{-3}-D_*^{-3})L(AW),\\
 P_-&=\frac{A^2}{CD_*^3}\{L(AW+\delta)-L(AW)\}.
\end{align*}
The identical decomposition with $L_T$ has the same nonnegative signs.
Mean-value estimates using $L_T'\le1/(M+z)$ give all the same majorants.
Thus the pairing remains uniform for the actual common cutoff, rather than
being an algebraic manipulation of separately divergent unregulated integrals.

\subsection{Four unit-cube sectors and their Jacobians}

There can also be joint faces $A,u\to0$ or $A,1-u\to0$. Split the inner
$u$ interval into two halves and parameterize them by
\[
 u=r/2\quad\hbox{or}\quad u=1-r/2,
 \qquad v=(1-u)z,\quad w=(1-u)(1-z).
\]
The inner Jacobian is $(1-u)/2$. Set
\[
 V=1-u,\qquad h_0=\frac{1-r/2}{2},\qquad
 H=rh_0,\qquad \frac14\le h_0\le\frac12.
\]
Also put $b_0=B/r$ by continuity; in the lower and upper halves respectively,
\[
 b_0=\tfrac12V^2z(1-z),\qquad
 b_0=\tfrac14ur\,z(1-z).
\]
In each half split the $(A,r)$ square into $A\le r$ and $r\le A$.
For $A\le r$, let $A=rt$; the Jacobian is $r$, and
\[
 D_*=rd,\quad d=t+(1-rt)h_0\ge\frac14,\quad
 N=tG+Cb_0s_1,\quad \delta=CN/d.
\]
Including the $rV/2$ Jacobian, the two positive chart densities are
\begin{align*}
 \widehat P_+&=\frac{rVh_0}{2}
       \frac{d^2+td+t^2}{d^3}\frac WM\ell(AW/M),\\
 \widehat P_-&=\frac{Vt^2N}{2d^4(M+AW)}
                  \ell\left(\frac{\delta}{M+AW}\right).
\end{align*}
For $r\le A$, let $r=At$; the Jacobian is $A$, and
\[
 D_*=Ad,\quad d=1+Ct h_0\ge1,\quad
 N=G+Ct b_0s_1,\quad\delta=CN/d.
\]
The corresponding positive densities are
\begin{align*}
 \widehat P_+&=\frac{AVt h_0}{2}
       \frac{d^2+d+1}{d^3}\frac WM\ell(AW/M),\\
 \widehat P_-&=\frac{VN}{2d^4(M+AW)}
                  \ell\left(\frac{\delta}{M+AW}\right).
\end{align*}
Finally, $a=x$, $b=(1-x)y$, $c=(1-x)(1-y)$ has outer Jacobian $1-x$.
The two inner halves and the two order sectors therefore give four
five-dimensional unit-cube integrals. All explicit radial poles are
cancelled, and the reduced denominators have the positive lower bounds
shown above. These are exact finite coordinate maps; no numerical accuracy
claim is needed for their derivation.

\subsection{An explicit linear parameter-norm bound}

Let $S=\max(s_1,s_2,s_3)$. The vector $k_0$ is a subconvex combination
of $p_1,-p_2$, and $k_0-p_1=-ap_1+bp_3$. Therefore their squared norms
are at most $S$, and
\[
 0\le W\le S/3,\qquad 0\le G\le S/3.
\]
Furthermore
\[
 HG-Bs_1=|Hk_0-uwp_1|^2\ge0,
 \qquad 0\le\delta\le CG.
\]
We bound $P_+$ and $P_-$ separately, which controls the absolute value of
their paired difference at every parameter point.

Since $D_*\ge\max(A,H)$, the stable positive term obeys $P_+\le3W/M$.
Using
$\int_{\Delta_o}W\,da\,db=(s_1+s_2+s_3)/24$ and
$\operatorname{vol}(\Delta_i)=1/2$, its integral is at most
$3S/(16M)$. For the negative term use
\[
 AG+CBs_1\le S\{AH+(A+Cu)vw\},\qquad
 \frac{A^3H}{D_*^4}\le1,
 \qquad (A+Cu)(1-u)\le D_*.
\]
The $AH$ contribution is at most $S/(4M)$ after integrating the two
fraction simplexes. The other contribution is at most $(S/(12M))I$,
where $1/12$ is the outer simplex volume $1/2$ times
$\int_0^1z(1-z)\,dz=1/6$, and
\[
 I=\int_0^1\int_0^1
 \frac{A^2(1-u)^2}{\{A+(1-A)u(1-u)\}^3}\,dA\,du.
\]
For $u\le1/2$, $A\le1/2$, use $D_*\ge A+u/4$ and extend the
$u$ integral to infinity. This part is at most
$\int_0^{1/2}2\,dA=1$. For $u\le1/2$, $A\ge1/2$, use $D_*\ge A$,
which gives $(7/24)\log2$. For $u\ge1/2$, put $v=1-u$. Since
\[
 D_*\ge\frac{A+v}{2},\qquad
 \frac{A^2v^2}{(A+v)^4}\le\frac1{16},
\]
the remaining integral is at most
$\int_0^1\int_0^{1/2}(A+v)/2\,dv\,dA=3/16$. Hence
\[
 I\le\frac{19}{16}+\frac7{24}\log2.
\]
Adding the three contributions gives the explicit parameter $L^1$ bound
\[
 \int|\mathcal F|\le C_N\frac SM,\qquad
 C_N=\frac{103}{192}+\frac{7\log2}{288}.
\]
The same pointwise nonnegative majorants hold for the common-cutoff paired
density. They prove parameter dominated convergence at all the individual
and joint UV faces, not merely at a chosen momentum point.

\subsection{A global logarithmic-square bound}

Set $T_0=S/(3M)$, $L=\log(1+T_0)$, and
\[
 f(T_0)=\int_0^{T_0}\frac{\log(1+t)}{t}\,dt
            \le L+\frac12L^2.
\]
For completeness the derivative of the difference on the right is
$(T_0-\log(1+T_0))/(T_0(1+T_0))\ge0$, and the difference vanishes at zero.
The first positive term is bounded by
\[
 P_+\le\frac3A\log(1+AT_0).
\]
The two simplex volumes are each $1/2$, so its integral is at most
$3f(T_0)/4\le3L/4+3L^2/8$.

For the second term use $\delta\le CS/3$ and
\[
 \int_{\Delta_i}\frac{A^2}{D_*^3}\,du\,dv
       \le\min\left(\frac1C,\frac1{2A}\right).
\]
The $1/(2A)$ bound follows immediately from $D_*\ge A$ and inner volume
$1/2$. For the $1/C$ bound, symmetrize the factor $1-u$ and use
$u(1-u)\ge u/2$ on $u\le1/2$; the resulting half-interval integral is
bounded by $\int_0^\infty A^2/(A+Cu/2)^3\,du=1/C$.
It follows that the second term integrates to at most
\[
 \frac12\int_0^1\frac{\log(1+CT_0)}C
                \min\left(\frac1C,\frac1{2A}\right)dA.
\]
Split at $A=1/2$. The first part is at most $L$, and the second at most
$f(T_0)/2$, giving $3L/2+L^2/4$. Altogether
\[
 |J_{\mathrm{tr}}|\le\int|\mathcal F|
 \le\min\left\{C_N\frac SM,\frac94L+\frac58L^2\right\}.
\]
This is compact-uniform when $S\le P^2$, $M\ge m_0^2>0$, and it has
at most logarithmic-square growth at a large common external scale. It is
a parameter amplitude bound, not a quadrature error or an original
momentum $L^1$ constant.

\subsection{Independent original-momentum absolute convergence}

Put $P=\max_i|p_i|$, $\Lambda=1+P/m$. On the physical ray
$p_1\mapsto tp_1$, $p_2\mapsto tp_2$, denote the factors by $a_t,h_t$.
The complete overall remainder is
\[
 RI=\int_0^1\{a_t'(h_t-D(q)^3)+a_t h_t'\}\,dt.
\]
The cancellation is made before taking absolute values. Use the uniform
propagator and first-derivative bounds of Section~\ref{sec:scheme}.

\paragraph{Hard inner loop.}
If $|q|\ge2(|k|+P+m)$, telescoping propagator differences gives
\begin{align*}
 |h_t-D(q)^3|&\le128(|k|+P)\langle q\rangle^{-7},\\
 |h_t'|&\le64P\langle q\rangle^{-7},\\
 |a_t|&\le\Lambda^4\langle k\rangle^{-6},\qquad
 |a_t'|\le4P\Lambda^5\langle k\rangle^{-7}.
\end{align*}
Thus
\[
 |RI|\le P(512\Lambda^6+64\Lambda^4)
                      \langle k\rangle^{-6}\langle q\rangle^{-7}.
\]
The hard $q$ tail is at most $1/(2\langle k\rangle)$, and
$\int_k\langle k\rangle^{-7}=8/(15m)$. This proves a finite norm in
the region that contains the logarithmically divergent inner triangle
before subtraction.

\paragraph{Hard outer loop.}
If $|k|\ge2(|q|+P+m)$, direct estimates on the first-order physical
remainder give
\[
 |RI|\le3072P\langle k\rangle^{-11}\langle q\rangle^{-2}
            +128P\langle k\rangle^{-7}\langle q\rangle^{-6}.
\]
The two hard tail integrals are bounded respectively by
$1/(160\langle q\rangle^5)$ and $1/(2\langle q\rangle)$.
The remaining $q$ integral in both cases is bounded by
$\int_q\langle q\rangle^{-7}=8/(15m)$.

\paragraph{Comparable loops, including a soft sum.}
On the remaining region let $K=\max(|k|,|q|)\ge K_0=8(P+m)$.
Then $|k|,|q|\ge K/4$, and the shifted outer momenta are at least
$K/8$. Retain $l=k+q$ as a possibly soft variable. Uniformly in $t$,
\begin{align*}
 |a_t'(h_t-D(q)^3)+a_t h_t'|
 \le{}&4P\,8^9K^{-9}D(l)D(l-tp_1)\\
          &+2P\,8^8K^{-8}D(l)\langle l-tp_1\rangle^{-3}
            +4P\,8^{13}K^{-13}.
\end{align*}
Integrating this inequality over $0\le t\le1$ bounds $|RI|$; the same
bounds below are uniform in $t$. On a dyadic shell $K\in[K_n,2K_n]$,
the $k$ ball volume is at most $64K_n^6/6$, and $|l|\le4K_n$.
Shift comparison and radial integration give
\begin{align*}
 \int_{|l|\le4K_n}D(l)D(l-tp_1)\frac{d^6l}{\pi^3}
        &\le8\Lambda^2K_n^2,\\
 \int_{|l|\le4K_n}D(l)\langle l-tp_1\rangle^{-3}
                                  \frac{d^6l}{\pi^3}
        &\le4\Lambda^3K_n.
\end{align*}
The purely hard last term uses product ball volume at most
$4096K_n^{12}/36$. Hence each shell is bounded by $P/K_n$ times the
explicit constant
\[
 \frac{1024}{3}8^9\Lambda^2
       +\frac{256}{3}8^8\Lambda^3+\frac{4096}{9}8^{13}.
\]
Since $\sum_nK_n^{-1}=2/K_0$, this region is absolutely integrable,
including $k\simeq-q$ where the two inner shifted lines are soft.

\paragraph{Compact remainder and uniform cutoff removal.}
For $K\le K_0$, positive mass gives $|RI|\le10Pm^{-13}$ and product
ball volume at most $K_0^{12}/36$. Combining the compact bound with the
three large-momentum regions proves absolute convergence in the original
twelve dimensions. This is an independent proof, not an inference from the
parameter integral.

The telescoping difference estimates and all first derivatives remain valid
for $D_\eta$, because its positive Schwinger moments are bounded by the
zero-cutoff moments. Thus these momentum majorants are uniform in $\eta$.
Momentum dominated convergence gives $\int RI_\eta\to\int RI$, while
the common-cutoff parameter pairing and the integrable majorants above give
the same left side tending to $J_{\mathrm{tr}}$. The two limits coincide.
If external momenta are bounded and $m\ge m_0>0$, use $m_0$ in the
majorants from the start to obtain compact-uniform control. The proven
result is this family of three massive Euclidean graphs, with its fixed
finite BPH scheme and physical leg mappings.
